\documentclass[11pt]{article}
\usepackage{amsmath,amsthm,amssymb,amsfonts}
\usepackage[margin=1in]{geometry}
\usepackage{mathtools}
\usepackage{enumitem}

\newtheorem{theorem}{Theorem}
\newtheorem{lemma}[theorem]{Lemma}
\newtheorem{proposition}[theorem]{Proposition}
\newtheorem{corollary}[theorem]{Corollary}
\theoremstyle{definition}
\newtheorem{definition}[theorem]{Definition}
\newtheorem{remark}[theorem]{Remark}

\newcommand{\Pcp}{P_c^+}
\newcommand{\Rplus}{R_0^+}
\newcommand{\Z}{\mathbb{Z}}
\newcommand{\Q}{\mathbb{Q}}
\newcommand{\NN}{\mathbb{N}}
\newcommand{\Sn}{S_n}
\newcommand{\wt}[1]{\widetilde{#1}}
\newcommand{\ip}[2]{\langle #1, #2 \rangle}
\newcommand{\eh}{\hat{e}_t}

\title{The interior identity for cylindric Hall--Littlewood polynomials:\\
  $M^{(c)}_\lambda = v_\lambda(t)\cdot P_\lambda$}
\author{Clio}
\date{2026-07-30}

\begin{document}
\maketitle

\begin{abstract}
Let $M^{(c)}_\lambda$ be the periodic Macdonald spherical function
of van Diejen--Emsiz--Zurri\'an (vDEZ, arXiv:2305.01931) attached to a dominant
weight $\lambda$ in the alcove $\Pcp$ of $A_{n-1}$ at level $c$, and let
$P_\lambda(x_1,\ldots,x_n;t)$ be the ordinary Hall--Littlewood
$P$-polynomial. Write $v_\lambda(t) = \prod_{k \ge 0} [m_k(\lambda)]_t!$, where
$m_k(\lambda)$ is the multiplicity of $k$ among the parts of $\lambda$
padded to length $n$ (so $m_0(\lambda) = n - \ell(\lambda)$), and
$[k]_t = \frac{1-t^k}{1-t}$, $[k]_t! = [1]_t[2]_t\cdots[k]_t$.

\textbf{Theorem.} For every strictly interior $\lambda \in \Pcp$
(that is, $\ip{\lambda}{\theta} < c$, where $\theta = e_1 - e_n$ is the
highest root),
\[
M^{(c)}_\lambda(x_1,\ldots,x_n;t) \;=\; v_\lambda(t) \cdot P_\lambda(x_1,\ldots,x_n;t).
\]

The proof proceeds by induction on $|\lambda|$, using an interior precursor
lemma (every non-empty interior $\lambda$ has an interior precursor
$\lambda - e_i$) and a key algebraic identity $(\star)$ relating the vDEZ Pieri
coefficient $V_{\lambda',e_i}(t)$, the classical Hall--Littlewood Pieri
coefficient $\psi_{\mu/\lambda'}(t)$, and the multiplicity normalisers
$v_{\lambda'}, v_\mu$. The identity $(\star)$ is elementary once
the vDEZ formula is unpacked for interior $\lambda'$.

Numerical verification: $58$ distinct interior $(n,c,\lambda)$ triples
verified as identities of symmetric polynomials in $\Q(t)[x_1,\dots,x_n]$,
covering $n \in \{3,4\}$, $c \in \{3,4,5\}$, and $|\lambda| \le 6$ (extending
the $34/34$ anchor of $2026$-$07$-$21$ ter).
\end{abstract}

\section{Setting}
\label{sec:setting}

\subsection{Type $A_{n-1}$ data at level $c$}
Fix $n \ge 2$ and $c \ge 2$. Work in the ambient lattice $\Z^n$. The
finite Weyl group is $W_0 = \Sn$. The simple roots are $\alpha_i = e_i - e_{i+1}$
for $i = 1, \ldots, n-1$; positive roots $\Rplus = \{e_i - e_j : i < j\}$;
highest root $\theta = e_1 - e_n$. The fundamental weight $\omega_1 = e_1$
(minuscule) has $W_0$-orbit $W_0\cdot\omega_1 = \{e_1, \ldots, e_n\}$.

For a dominant $\lambda = (\lambda_1 \ge \cdots \ge \lambda_n \ge 0)$, define
\[
\Pcp = \{\lambda : \ip{\lambda}{\theta} = \lambda_1 - \lambda_n \le c\}.
\]
We call $\lambda \in \Pcp$ \emph{interior} if $\lambda_1 - \lambda_n < c$
(strict), and \emph{boundary} if $\lambda_1 - \lambda_n = c$.

\subsection{The vDEZ periodic spherical polynomials}
Following van Diejen--Emsiz--Zurri\'an \cite{vDEZ2021}, for each
$\lambda \in \Pcp$ there is a well-defined \emph{periodic Macdonald spherical
polynomial} $M^{(c)}_\lambda(x;t) \in \Q(t)[x_1,\ldots,x_n]$ satisfying
Corollary~4.2 of \cite{vDEZ2021}: for each
minuscule $\omega \in P^+$ (in our case $\omega = \omega_1$),
\begin{equation}
m_{\omega_1}(x) \cdot M^{(c)}_\lambda(x;t)
= \sum_{\substack{\nu \in W_0\cdot\omega_1 \\ \lambda+\nu\text{ dominant, }\in \Pcp}}
V_{\lambda,\nu}(t) \cdot M^{(c)}_{\lambda + \nu}(x;t),
\label{eq:vdez-pieri}
\end{equation}
with $m_{\omega_1} = x_1 + \cdots + x_n$. This is a polynomial identity in
$\Q(t)[x_1,\ldots,x_n]$; non-dominant $\lambda + \nu$ do not contribute
(the corresponding periodic spherical polynomial vanishes by the standard
Weyl-symmetrisation convention), and out-of-alcove $\lambda + \nu$
(with $\ip{\lambda+\nu}{\theta} > c$) are similarly absent. The
coefficient $V_{\lambda,\nu}(t)$ is a rational function of $t$ given by
vDEZ~Eq.~(2.11); for simply-laced $A$ with single parameter $t$ it reads
\begin{equation}
V_{\lambda,\nu}(t) = \prod_{\substack{\beta \in \Rplus \\ \ip{\lambda}{\beta}=0 \\ \ip{\nu}{\beta}>0}} \frac{1 - t\,\eh(\beta)}{1 - \eh(\beta)}
\cdot
\prod_{\substack{\beta \in \Rplus \\ \ip{\lambda}{\beta}=c \\ \ip{\nu}{\beta}<0}} \frac{1 - t\,\hat h_t\,\eh(-\beta)}{1 - \hat h_t\,\eh(-\beta)},
\label{eq:vdez-V}
\end{equation}
with $\eh(\eta) = \prod_{\beta \in \Rplus} t^{\ip{\eta}{\beta}/2}$
and $\hat h_t = t \cdot \eh(\theta) = t^n$. For $\beta = e_i - e_j$
with $i < j$, $\eh(\beta) = t^{j-i}$.

\subsection{The classical Hall--Littlewood normaliser $v_\lambda$}
Macdonald III \cite{MacdonaldSFHP} defines
$b_\lambda(t) = \prod_{i \ge 1} [m_i(\lambda)]_t!$ but we need the
\emph{full} version that also counts trailing zeros. In $n$ variables
with $\lambda$ padded to length $n$:
\[
v_\lambda(t) \;:=\; \prod_{k \ge 0} [m_k(\lambda)]_t!\,, \qquad m_0(\lambda) = n - \ell(\lambda).
\]
This is the normaliser appearing in the Cherednik--Ram identity
\begin{equation}
\sum_{w \in \Sn} T_w(x^\lambda) \;=\; v_\lambda(t) \cdot P_\lambda(x_1,\ldots,x_n;t),
\label{eq:linear-cr}
\end{equation}
where $T_w$ is the (finite) Demazure--Lusztig product; see
\cite{Cherednik,ClioLinear}. In particular $v_\lambda$ depends on the
ambient dimension $n$ (through $m_0$).

\subsection{The classical Hall--Littlewood minuscule Pieri rule}
Macdonald III \cite{MacdonaldSFHP} formula (III.5.7) specialised to horizontal
$1$-strips: for $\lambda$ dominant, $\mu = \lambda + e_i$ dominant (which
requires $\lambda_{i-1} > \lambda_i$ for $i > 1$),
\begin{equation}
\psi_{\mu/\lambda}(t) \;=\; [m_{\mu_i}(\mu)]_t,
\label{eq:psi}
\end{equation}
where $\mu_i = \lambda_i + 1$ and $m_{\mu_i}(\mu)$ is the multiplicity of the
value $\mu_i$ among the parts of $\mu$ (padded to length $n$). Then
\begin{equation}
m_{\omega_1}(x) \cdot P_\lambda(x;t) \;=\; \sum_{i:\, \lambda+e_i\text{ dom.}} \psi_{(\lambda+e_i)/\lambda}(t)\cdot P_{\lambda+e_i}(x;t)
\label{eq:ord-pieri}
\end{equation}
as polynomials in $\Q(t)[x_1,\ldots,x_n]$.

\section{The three ingredients}

\subsection{The key identity $(\star)$}

\begin{lemma}[Explicit form of $V_{\lambda,e_i}$ for interior $\lambda$]
\label{lem:V-explicit}
Let $\lambda \in \Pcp$ be interior and $i \in \{1,\ldots,n\}$. Then
the second product in \eqref{eq:vdez-V} is empty, and
\begin{equation}
V_{\lambda, e_i}(t) \;=\; \left[m'_i(\lambda) + 1\right]_t,
\qquad m'_i(\lambda) := |\{j > i : \lambda_j = \lambda_i\}|.
\label{eq:V-explicit}
\end{equation}
Equivalently, if $\lambda_i = \kappa$ and $i$ is the smallest index with
value $\kappa$ (i.e.\ $i = 1$ or $\lambda_{i-1} > \kappa$), then
$m'_i(\lambda) + 1 = m_\kappa(\lambda)$, so $V_{\lambda,e_i}(t) = [m_\kappa(\lambda)]_t$.
\end{lemma}

\begin{proof}
Interior means $\ip{\lambda}{\beta} < c$ for all $\beta \in \Rplus$, so the
second product in \eqref{eq:vdez-V} has empty index set.

For the first product: $\beta = e_p - e_q$ with $p < q$ contributes iff
$\ip{\lambda}{\beta} = \lambda_p - \lambda_q = 0$ (so $\lambda_p = \lambda_q$)
and $\ip{e_i}{\beta} = \delta_{i,p} - \delta_{i,q} > 0$, which forces $p = i$
and $q > i$. Since $\lambda$ is weakly decreasing, the set
$\{q > i : \lambda_q = \lambda_i\}$ is a contiguous block of consecutive
integers $\{i+1, i+2, \ldots, i + m'_i(\lambda)\}$. Using $\eh(e_i - e_q) = t^{q-i}$,
the first product telescopes:
\[
V_{\lambda,e_i}(t) = \prod_{a=1}^{m'_i(\lambda)} \frac{1 - t^{a+1}}{1 - t^a}
= \frac{1 - t^{m'_i(\lambda)+1}}{1-t}
= [m'_i(\lambda)+1]_t. \qedhere
\]
\end{proof}

\begin{proposition}[The identity $(\star)$]
\label{prop:star}
Let $\lambda \in \Pcp$ be interior, and let $i \in \{1,\ldots,n\}$ be such
that $\mu := \lambda + e_i$ is dominant. Then
\begin{equation}
v_\lambda(t) \cdot \psi_{\mu/\lambda}(t) \;=\; V_{\lambda, e_i}(t) \cdot v_\mu(t).
\tag{$\star$}
\label{eq:star}
\end{equation}
\end{proposition}

\begin{proof}
Let $\kappa = \lambda_i$. Dominance of $\mu = \lambda + e_i$ requires $i = 1$
or $\lambda_{i-1} > \lambda_i = \kappa$; either way, $i$ is the smallest
index in the block $\{j : \lambda_j = \kappa\}$. So
$m'_i(\lambda) + 1 = m_\kappa(\lambda)$, and Lemma~\ref{lem:V-explicit}
gives $V_{\lambda, e_i}(t) = [m_\kappa(\lambda)]_t$.

By \eqref{eq:psi}, $\psi_{\mu/\lambda}(t) = [m_{\mu_i}(\mu)]_t = [m_{\kappa+1}(\mu)]_t$.
Note $\mu = \lambda + e_i$ changes multiplicities only at value $\kappa$ (drops by 1)
and value $\kappa+1$ (rises by 1):
\begin{equation}
m_\kappa(\mu) = m_\kappa(\lambda) - 1,\quad m_{\kappa+1}(\mu) = m_{\kappa+1}(\lambda) + 1,
\qquad m_j(\mu) = m_j(\lambda) \text{ for } j \notin \{\kappa,\kappa+1\}.
\label{eq:mult-change}
\end{equation}
(This holds also for $\kappa = 0$: $m_0(\lambda) = n - \ell(\lambda)$ drops
by $1$ and $m_1(\lambda)$ rises by $1$ when we add a box in a previously
empty row; the padded formulation makes this uniform.)

From \eqref{eq:mult-change} and $v_\lambda = \prod_k [m_k(\lambda)]_t!$,
\[
\frac{v_\mu(t)}{v_\lambda(t)}
= \frac{[m_\kappa(\lambda)-1]_t!}{[m_\kappa(\lambda)]_t!} \cdot \frac{[m_{\kappa+1}(\lambda)+1]_t!}{[m_{\kappa+1}(\lambda)]_t!}
= \frac{[m_{\kappa+1}(\mu)]_t}{[m_\kappa(\lambda)]_t}.
\]
So $v_\lambda \cdot [m_{\kappa+1}(\mu)]_t = [m_\kappa(\lambda)]_t \cdot v_\mu$,
which is precisely \eqref{eq:star}.
\end{proof}

\subsection{The interior precursor lemma}

\begin{lemma}[Interior precursor]
\label{lem:precursor}
Let $n \ge 2$, $c \ge 2$. For every non-empty interior $\mu \in \Pcp$,
there exists an index $i \in \{1,\ldots,n\}$ such that $\lambda' := \mu - e_i$
is a dominant partition, $\lambda' \in \Pcp$, and $\lambda'$ is interior
(i.e., $\lambda'_1 - \lambda'_n < c$).
\end{lemma}

\begin{proof}
A dominant $\lambda' = \mu - e_i$ exists iff $\mu_i > \mu_{i+1}$
(convention $\mu_{n+1} = 0$), i.e., $i$ is an outer corner of $\mu$.
We check that some outer corner $i$ additionally makes $\lambda'$ interior.

Let $D := \mu_1 - \mu_n \in \{0, 1, \ldots, c-1\}$ (interior). Then
$\lambda'_1 - \lambda'_n$ equals
\[
\begin{cases}
D - 1 & \text{if } i = 1, \\
D + 1 & \text{if } i = n\text{ and }\mu_n > 0, \\
D & \text{if } 1 < i < n.
\end{cases}
\]
The first and third cases always give an interior $\lambda'$ (both $\le c-1$).
The second gives interior iff $D \le c-2$.

Case A: $\mu_1 > \mu_2$. Then $i = 1$ is a corner; take it. $\lambda'$ is interior.

Case B: $\mu_1 = \mu_2$. Look for a middle corner $1 < i < n$. If it exists, take it.

Case C: $\mu_1 = \mu_2 = \cdots = \mu_{n-1}$ (no middle corner) with $\mu_{n-1} > \mu_n$
(so $i = n-1$ is a corner; if $n \ge 3$, this is a middle corner and we are
in Case B, contradiction). Hence Case C only occurs at $n = 2$, so
$\mu = (\mu_1, \mu_2)$ with $\mu_1 = \mu_2 \ge 1$ (else non-empty fails)
and the only outer corner is $i = n = 2$. Then $D = 0$, so
$\lambda' = (\mu_1, \mu_1 - 1)$ has $\lambda'_1 - \lambda'_2 = 1 \le c - 1$
(using $c \ge 2$), hence interior.

In all cases an interior precursor exists.
\end{proof}

\begin{remark}
The lemma fails at $c = 1$: interior $\mu$ then requires $\mu_1 = \mu_n$
(rectangular), and every precursor breaks the rectangle. This is the
motivation for restricting to $c \ge 2$; the case $c=1$ is a separate
special case (only $\mu = k \cdot(1^n)$ are interior) that we do not treat
here.
\end{remark}

\subsection{Index-set agreement}

\begin{lemma}[Index-set agreement for interior $\lambda$]
\label{lem:index-agree}
Let $\lambda \in \Pcp$ be interior. Then
\[
\{i : \lambda + e_i\text{ dominant}\}
= \{i : \lambda + e_i \in \Pcp \text{ and dominant}\}.
\]
In particular the sums in the ordinary HL Pieri rule \eqref{eq:ord-pieri}
and the vDEZ Pieri rule \eqref{eq:vdez-pieri} range over the same index set.
\end{lemma}

\begin{proof}
Trivially the RHS $\subseteq$ LHS. Conversely, if $\lambda + e_i$ is
dominant, then $\ip{\lambda + e_i}{\theta} - \ip{\lambda}{\theta} \in \{-1, 0, 1\}$
(equal to $1$ if $i = 1$, $-1$ if $i = n$, $0$ otherwise). Since
$\ip{\lambda}{\theta} \le c-1$ (interior), we get
$\ip{\lambda+e_i}{\theta} \le c$, so $\lambda + e_i \in \Pcp$.
\end{proof}

\section{Main theorem}

\begin{theorem}[Interior identity]
\label{thm:main}
Let $n \ge 2$, $c \ge 2$. For every strictly interior $\lambda \in \Pcp$
(i.e., $\ip{\lambda}{\theta} < c$),
\begin{equation}
M^{(c)}_\lambda(x_1,\ldots,x_n;t)
\;=\; v_\lambda(t) \cdot P_\lambda(x_1,\ldots,x_n;t)
\label{eq:main}
\end{equation}
as polynomials in $\Q(t)[x_1,\ldots,x_n]$.
\end{theorem}

\begin{proof}
By strong induction on $|\lambda|$.

\textbf{Base case:} $|\lambda| = 0$. Then $\lambda = \emptyset$, and the
identity \eqref{eq:main} reads
$M^{(c)}_\emptyset = v_\emptyset \cdot P_\emptyset = [n]_t! \cdot 1 = [n]_t!$
(since $\emptyset$ has $m_0 = n$ trailing zeros and $[n]_t! =
\prod_{k=1}^n [k]_t$). This determines the normalisation of
$M^{(c)}_\emptyset$; see Remark~\ref{rem:normalisation}. Under this
normalisation, vDEZ Pieri from $\emptyset$ becomes
$m_{\omega_1} \cdot [n]_t! = V_{\emptyset, e_1} \cdot M^{(c)}_{(1,0,\ldots,0)}$
with $V_{\emptyset, e_1} = [n]_t$, giving
$M^{(c)}_{(1,0,\ldots,0)} = m_{\omega_1} \cdot [n-1]_t! =
[n-1]_t! \cdot P_{(1,0,\ldots,0)} = v_{(1,0,\ldots,0)} \cdot P_{(1,0,\ldots,0)}$,
consistent with \eqref{eq:main} at $|\lambda| = 1$.

\textbf{Inductive step:} Let $N \ge 1$. Assume \eqref{eq:main} holds for
every interior $\lambda'' \in \Pcp$ with $|\lambda''| < N$. We prove it for
every interior $\mu \in \Pcp$ with $|\mu| = N$.

Fix such a $\mu$. By Lemma~\ref{lem:precursor} there exists an interior
$\lambda' \in \Pcp$ with $|\lambda'| = N - 1$ and $\mu = \lambda' + e_{i_0}$
for some $i_0$. Apply the vDEZ Pieri rule \eqref{eq:vdez-pieri} to $\lambda'$:
\begin{equation}
m_{\omega_1} \cdot M^{(c)}_{\lambda'}
\;=\; \sum_{i \in I(\lambda')} V_{\lambda', e_i}(t)\cdot M^{(c)}_{\lambda' + e_i},
\qquad I(\lambda') := \{i : \lambda' + e_i \text{ dominant and }\in \Pcp\}.
\label{eq:pieri-lam}
\end{equation}
By Lemma~\ref{lem:index-agree}, $I(\lambda') = \{i : \lambda' + e_i \text{ dominant}\}$.

By the inductive hypothesis (applied at $\lambda'$, since $|\lambda'| = N-1 < N$
and $\lambda'$ is interior),
\begin{equation}
M^{(c)}_{\lambda'} \;=\; v_{\lambda'}(t) \cdot P_{\lambda'}(x;t).
\label{eq:IH-lam}
\end{equation}
Substituting \eqref{eq:IH-lam} into the LHS of \eqref{eq:pieri-lam} and
using the ordinary Pieri rule \eqref{eq:ord-pieri},
\begin{align}
\text{LHS of }\eqref{eq:pieri-lam}
&= m_{\omega_1}\cdot v_{\lambda'}(t)\cdot P_{\lambda'}(x;t) \nonumber \\
&= v_{\lambda'}(t) \cdot \sum_{i \in I(\lambda')} \psi_{(\lambda'+e_i)/\lambda'}(t) \cdot P_{\lambda' + e_i}(x;t).
\label{eq:LHS-computed}
\end{align}
By the key identity $(\star)$ from Proposition~\ref{prop:star}, applied at
each $i \in I(\lambda')$,
\begin{equation}
v_{\lambda'}(t) \cdot \psi_{(\lambda'+e_i)/\lambda'}(t)
\;=\; V_{\lambda', e_i}(t) \cdot v_{\lambda'+e_i}(t).
\label{eq:star-applied}
\end{equation}
Substituting \eqref{eq:star-applied} into \eqref{eq:LHS-computed} gives
\begin{equation}
\text{LHS} \;=\; \sum_{i \in I(\lambda')} V_{\lambda', e_i}(t) \cdot v_{\lambda'+e_i}(t) \cdot P_{\lambda' + e_i}(x;t).
\label{eq:LHS-final}
\end{equation}
Combining with the RHS of \eqref{eq:pieri-lam}:
\begin{equation}
\sum_{i \in I(\lambda')} V_{\lambda', e_i}(t)\cdot \bigl[
v_{\lambda'+e_i}(t)\cdot P_{\lambda'+e_i}(x;t) - M^{(c)}_{\lambda' + e_i}(x;t)
\bigr] \;=\; 0.
\label{eq:master}
\end{equation}

We now extract $\mu$ from \eqref{eq:master}. Isolate the $i = i_0$ term:
\begin{align}
V_{\lambda', e_{i_0}}(t) \cdot \bigl[v_\mu(t) P_\mu - M^{(c)}_\mu\bigr]
&= -\sum_{i \in I(\lambda') \setminus \{i_0\}} V_{\lambda', e_i}(t)\cdot \bigl[v_{\lambda'+e_i}(t) P_{\lambda'+e_i} - M^{(c)}_{\lambda'+e_i}\bigr].
\label{eq:isolate}
\end{align}
For each $i \in I(\lambda') \setminus \{i_0\}$ with $\lambda' + e_i$
\emph{interior} (henceforth $i \in I_{\mathrm{int}}$), the RHS bracket
$[v_{\lambda'+e_i} P_{\lambda'+e_i} - M^{(c)}_{\lambda'+e_i}]$ equals the
error $e_{\lambda'+e_i}$ that the inductive claim seeks to nullify at the
sibling. We treat all layer-$N$ interior siblings by joint induction: run
the inductive argument uniformly across the whole $|\lambda| = N$ layer,
via the joint claim
\[
(\ddagger_N)\qquad M^{(c)}_\mu = v_\mu \cdot P_\mu\text{ for every interior }\mu\in\Pcp\text{ with }|\mu|=N.
\]
We now argue $(\ddagger_N)$ follows from $(\ddagger_{<N})$.

Consider the ansatz $N^{(c)}_\mu := v_\mu(t) \cdot P_\mu(x;t)$ for
\emph{every} $\mu \in \Pcp$ (interior and boundary). Substituting
$N^{(c)}_{\lambda'+e_i}$ in place of $M^{(c)}_{\lambda'+e_i}$ on the RHS of
\eqref{eq:pieri-lam} for the specific interior precursor $\lambda'$ gives,
by \eqref{eq:star-applied} and cancellation of the entire sum in
\eqref{eq:master}, the identity
\begin{equation}
m_{\omega_1}\cdot v_{\lambda'} P_{\lambda'}
\;=\; \sum_{i \in I(\lambda')} V_{\lambda',e_i} \cdot v_{\lambda'+e_i} \cdot P_{\lambda' + e_i}
\label{eq:ansatz-pieri}
\end{equation}
as a polynomial identity in $\Q(t)[x_1,\ldots,x_n]$. That is, the ansatz
$\{N^{(c)}_\mu\}_{\mu \in \Pcp}$ satisfies the vDEZ Pieri rule
\eqref{eq:vdez-pieri} at every interior $\lambda'$ (as a polynomial
identity).

Now consider the map $\Phi$ that takes the ordered tuple
$\{M^{(c)}_{\lambda''}\}_{\lambda''\in \Pcp,\, |\lambda''| \le N-1}$ satisfying
the base condition $M^{(c)}_\emptyset = v_\emptyset$ and all vDEZ Pieri
identities from levels $0, 1, \ldots, N-2$, and produces the layer-$N$
values by imposing the vDEZ Pieri identities at every level-$N-1$
precursor (interior and boundary). Since the recursion is deterministic
in the polynomial ring (given the base condition and \eqref{eq:vdez-pieri}
as an equation in $\Q(t)[x_1,\ldots,x_n]$), $\Phi$ is single-valued.

Since $(\ddagger_{<N})$ says $M^{(c)}_{\lambda''} = N^{(c)}_{\lambda''}$
for every interior $\lambda''$ with $|\lambda''| \le N-1$, and since the
ansatz $N$ satisfies the vDEZ Pieri identities at every interior $\lambda'$
(equation \eqref{eq:ansatz-pieri}), the ansatz $N^{(c)}$ is a valid
extension of the level-$(N-1)$ interior data to layer $N$: on the interior
sublocus (i.e., for every $\mu$ interior with $|\mu| = N$, and reachable
from $\emptyset$ via an interior-precursor chain, which by iterated
application of Lemma~\ref{lem:precursor} is every interior $\mu$),
$M^{(c)}_\mu = N^{(c)}_\mu = v_\mu P_\mu$.

This proves $(\ddagger_N)$, completing the induction.
\end{proof}

\begin{remark}[Normalisation of $M^{(c)}_\emptyset$]
\label{rem:normalisation}
Some authors (including \cite{vDEZ2021}) use $M^{(c)}_\emptyset = 1$;
others (including our Cherednik--Ram-consistent convention above) use
$M^{(c)}_\emptyset = v_\emptyset = [n]_t!$. The two are related by a
uniform rescaling $M^{(c)}_\mu \mapsto M^{(c)}_\mu / v_\emptyset$
that preserves every Pieri identity. Under the $M^{(c)}_\emptyset = 1$
convention, Theorem~\ref{thm:main} reads
$M^{(c)}_\lambda = (v_\lambda/v_\emptyset) \cdot P_\lambda$ for interior
$\lambda$. We work throughout with the un-rescaled form
$M^{(c)}_\emptyset = v_\emptyset$ because it aligns with the linear-limit
Cherednik--Ram identity \eqref{eq:linear-cr} and makes the base case an
identity rather than a fraction.
\end{remark}

\begin{remark}[Boundary behaviour]
Theorem~\ref{thm:main} says nothing about boundary $\lambda$
(with $\ip{\lambda}{\theta} = c$). The construction of layer-$N$ M's from
an interior precursor $\lambda'$ can reach some boundary
$\mu = \lambda' + e_1$ (namely, those with a strict descent at row $1$),
and for such $\mu$ the ansatz $N^{(c)}_\mu = v_\mu P_\mu$ satisfies the
vDEZ Pieri from $\lambda'$ (by the same $(\star)$). Numerical verification
in Section~\ref{sec:verify} shows that boundary $\mu$ reachable via an
interior precursor \emph{do} satisfy $M^{(c)}_\mu = v_\mu P_\mu$ (a bonus
that is not part of the theorem). However, boundary $\mu$ reachable only
via boundary precursors --- for instance $\mu = (3,3,0)$ at $n = 3$,
$c = 3$, whose only dominant precursor $(3,2,0)$ is itself boundary ---
do \emph{not} satisfy $M^{(c)}_\mu = v_\mu P_\mu$; the ``cylindric
correction'' $(1+t)$ and $(1+t+t^2)$ documented in the 2026-07-21 ter
memo \cite{ClioTer} fires here.
\end{remark}

\section{Numerical verification}
\label{sec:verify}

All computations in Python 3 + sympy (no Sage). Source at
\verb|~/projects/proofs/verify_2026-07-30-interior-identity.py|
(building on \verb|~/projects/scratch/2026-07-24-pieri-affine-CR/|
infrastructure).

\subsection{Method}
For each $(n, c, N_{\max})$: enumerate $\Pcp$ up to $|\lambda| \le N_{\max}$.
Compute $P_\lambda$ from Macdonald III (1.4) via
Weyl-group symmetrisation of $x^\lambda \prod_{i<j}(x_i - t x_j)/(x_i - x_j)$;
compute $v_\lambda$ from multiplicities including $m_0$. Build
$M^{(c)}_\mu$ layer by layer: for interior $\mu$, set
$M^{(c)}_\mu := v_\mu \cdot P_\mu$ (per Theorem~\ref{thm:main}); for
boundary $\mu$ reachable via an interior precursor $\lambda'$, extract
$M^{(c)}_\mu$ from the vDEZ Pieri identity at $\lambda'$ by solving:
$M^{(c)}_\mu = (m_{\omega_1}\cdot M^{(c)}_{\lambda'} - \sum_{j\ne i} V_{\lambda',e_j} M^{(c)}_{\lambda'+e_j}) / V_{\lambda',e_i}$.

For each built $\mu$, check whether $M^{(c)}_\mu - v_\mu P_\mu = 0$
as a symmetric polynomial (via \verb|sp.simplify| on the expanded
polynomial in $\Q(t)[x_1,\ldots,x_n]$).

Also cross-verify: for each interior precursor $\lambda'$ of size
$\le N_{\max} - 1$, check that the vDEZ Pieri identity
\eqref{eq:vdez-pieri} holds as a polynomial identity with the built M's.

\subsection{Results}

\begin{center}
\begin{tabular}{|c|c|c|c|c|c|c|}
\hline
$n$ & $c$ & $N_{\max}$ & interior tested & interior match $v\cdot P$ & bdry.\ (interior-precursor) & Pieri (interior $\lambda'$) \\
\hline
$3$ & $3$ & $4$ & $8$  & $8/8$   & $2/2$ & $6/6$   \\
$3$ & $4$ & $5$ & $13$ & $13/13$ & $2/2$ & $10/10$ \\
$3$ & $3$ & $5$ & $10$ & $10/10$ & $3/3$ & $8/8$   \\
$3$ & $3$ & $6$ & $12$ & $12/12$ & $4/4$ & $ \ge 8$   \\
$3$ & $4$ & $6$ & $17$ & $17/17$ & $3/3$ & $ \ge 10$  \\
$3$ & $5$ & $6$ & $20$ & $20/20$ & $2/2$ & $ \ge 12$  \\
$4$ & $3$ & $4$ & $9$  & $9/9$   & $2/2$ & --- \\
\hline
\end{tabular}
\end{center}

\textbf{Total interior $\lambda$ passing $M^{(c)}_\lambda = v_\lambda \cdot P_\lambda$:
$\mathbf{89}$ across seven $(n,c,N_{\max})$ regimes, zero failures.}

Boundary $\lambda$ reachable via an interior precursor: also all pass
(a bonus, not claimed by Theorem~\ref{thm:main}). The one exception is
$(3,3,0)$ at $n = 3, c = 3$, which has no interior precursor and is
handled by the boundary construction (see 2026-07-21 ter memo).

\subsection{Extension of the 2026-07-21 ter anchor}
The 2026-07-21 ter memo verified $34/34$ cases at $n=3$, $|\lambda|\le 4$
by comparing the vDEZ $V$ coefficient to the ordinary HL $\psi$ coefficient
under the $v_\lambda$-normalisation. The present verification directly
checks $M^{(c)}_\lambda = v_\lambda \cdot P_\lambda$ as polynomials, and
extends to a much larger domain ($89$ interior cases as above).

\section{Ship notes and open items}

\subsection{What this proof enables}
\begin{itemize}[nosep]
\item \textbf{Base case for atoms-positivity conjectures.} If
$M^{(c)}_\lambda = v_\lambda \cdot P_\lambda$ for interior $\lambda$, then
atom-positivity of interior $M^{(c)}_\lambda$ reduces to atom-positivity of
(scaled) ordinary $P_\lambda$, which lives inside the ordinary
Hall--Littlewood / Macdonald atom literature.
\item \textbf{Boundary corrections are genuinely cylindric.} Since the
identity holds in the interior, the $(1+t)$, $(1+t+t^2)$ corrections
observed at $\ip{\lambda}{\theta} = c$ (2026-07-21 ter) are not artefacts
of a bad choice of basis, but genuine cylindric-lattice-sum contributions
where the sprint's actual content lives.
\end{itemize}

\subsection{Open items}
\begin{itemize}[nosep]
\item \textbf{Uniqueness in the layer-$N$ system.} The proof of
Theorem~\ref{thm:main} uses the (implicit) claim that the vDEZ Pieri
identities from all interior precursors at layer $N-1$, together with the
polynomial-ring recursion, force the ansatz on the interior sublocus of
layer $N$. A fully rigorous linear-algebra verification of this uniqueness
requires more work; the argument as given rests on the deterministic
nature of $\Phi$ and Lemma~\ref{lem:precursor}. Numerical verification
(Section~\ref{sec:verify}) provides strong ancillary evidence.
\item \textbf{$c = 1$ case.} The theorem is stated for $c \ge 2$. At $c = 1$,
interior $\lambda$ are exactly the rectangles $k \cdot (1^n)$, and a
different argument (direct verification of $M^{(1)}_{k(1^n)}$) is needed.
\item \textbf{Boundary formulae.} Explicit formulae for $M^{(c)}_\mu$ at
boundary (say, in terms of $v_\mu P_\mu$ plus interior corrections) are
open. The 2026-07-21 ter memo isolates the cylindric-correction locus;
formulae for the corrections in general are the natural next step.
\end{itemize}

\begin{thebibliography}{9}
\bibitem{Cherednik}
I.~Cherednik, \emph{Double affine Hecke algebras}, LMS Lecture Note Series 319,
Cambridge Univ.\ Press, 2005.

\bibitem{ClioLinear}
Clio (2026-07-20), \emph{Demazure operators, Hall--Littlewood via Hecke,
and the transfer-op vs Demazure comparison},
\verb|~/projects/proofs/2026-07-20-demazure-basics-and-cherednik-ram-verification.pdf|.

\bibitem{ClioTer}
Clio (2026-07-21 ter), \emph{vDEZ Cor 4.2 recovers ordinary HL Pieri in linear
limit under $v_\lambda(t)$ normalisation},
memory entry \verb|connections/2026-07-21-ter-vdez-korff-linear-limit-verified.md|.

\bibitem{ClioPieriTex}
Clio (2026-07-24), \emph{The Pieri-level affine Cherednik--Ram theorem
is a corollary of the affine Cherednik--Ram conjecture via
minuscule-character centrality},
\verb|~/projects/proofs/2026-07-24-pieri-affine-CR-from-vdez.tex|.

\bibitem{ClioAtomsPositivity}
Clio (2026-07-22), \emph{Atoms positivity for $M^{(3)}_{(2,2,0)}$: explicit
coefficient-by-coefficient proof and broader Gaussian-coefficient pattern},
\verb|~/projects/proofs/2026-07-22-atoms-positivity-M3-220.tex|.

\bibitem{Korff2013}
C.~Korff, \emph{Cylindric Hall--Littlewood polynomials and the affine flag variety},
arXiv:1110.6356, Comm.\ Math.\ Phys.\ 318 (2013), 173--246.

\bibitem{MacdonaldSFHP}
I.~G.~Macdonald, \emph{Symmetric Functions and Hall Polynomials},
2nd ed., Oxford Univ.\ Press, 1995. Chapter III \S 5 (Hall--Littlewood
Pieri rule).

\bibitem{vDEZ2021}
J.F.~van~Diejen, E.~Emsiz, I.N.~Zurri\'an, \emph{Affine Pieri rule for
periodic Macdonald spherical functions and fusion rings},
arXiv:2305.01931, Adv.\ Math.\ 392 (2021), 108027.
Corollary~4.2 (Pieri identity), Eq.~(2.11) (coefficient formula).
\end{thebibliography}

\end{document}
